\documentclass[11pt, letterpaper]{article}

% --- Idioma y codificación ---
\usepackage[spanish]{babel}
\usepackage[utf8]{inputenc}
\usepackage[T1]{fontenc}
\usepackage{lmodern}

% --- Matemáticas ---
\usepackage{amsmath, amssymb, amsthm}

% --- Formato ---
\usepackage[margin=2.5cm]{geometry}
\usepackage{parskip}
\usepackage{booktabs}
\usepackage{hyperref}
\hypersetup{colorlinks=true, linkcolor=blue!70!black, urlcolor=blue!70!black}

% --- Entornos de teoremas ---
\theoremstyle{definition}
\newtheorem{definition}{Definición}[section]
\theoremstyle{plain}
\newtheorem{theorem}{Teorema}[section]
\newtheorem{lemma}[theorem]{Lema}
\newtheorem{corollary}[theorem]{Corolario}
\theoremstyle{remark}
\newtheorem{remark}{Observación}[section]

\title{%
  T01 --- Demostración formal:\\[4pt]
  Costos de búsqueda en encadenamiento separado
}
\author{Curso Linux--C--Rust --- B15 Estructuras de datos en C}
\date{}

\begin{document}
\maketitle
\tableofcontents

% ===================================================================
\section{Modelo}
% ===================================================================

\begin{definition}
Array $T[0..m{-}1]$ con encadenamiento separado: cada bucket $T[j]$
contiene una lista enlazada. Nuevos elementos se insertan al
\textbf{inicio} de la lista (prepend). Bajo SUHA:
$\Pr[h(k) = j] = 1/m$ para toda clave~$k$, independientemente de las
demás. Notación: $n$ elementos, $m$ buckets,
$\alpha = n/m$, $L_j$ = longitud de la cadena en~$T[j]$.
Del T03: $E[L_j] = \alpha$ para todo~$j$.
\end{definition}

% ===================================================================
\section{Demostración 1: búsqueda fallida
  \texorpdfstring{$\Theta(1 + \alpha)$}{Theta(1 + alpha)}}
% ===================================================================

\begin{theorem}\label{thm:miss}
El número esperado de comparaciones en una búsqueda fallida es
exactamente $\alpha$.
\end{theorem}

\begin{proof}
Sea $k$ una clave que no está en la tabla. La búsqueda calcula
$j = h(k)$ en $O(1)$ y recorre \textbf{toda} la lista en $T[j]$,
comparando $k$ con cada elemento sin encontrarlo. El número de
comparaciones es exactamente~$L_j$.

Bajo SUHA, $h(k)$ es uniformemente aleatorio e independiente de los
$n$ elementos ya insertados:
\[
  E[\text{comparaciones}]
  = E[L_{h(k)}]
  = \sum_{j=0}^{m-1} \Pr[h(k) = j] \cdot E[L_j]
  = \sum_{j=0}^{m-1} \frac{1}{m} \cdot \alpha
  = \alpha
\]

Incluyendo el costo $O(1)$ de calcular el hash y acceder al bucket:
\[
  T_{\text{miss}} = \Theta(1 + \alpha) \qedhere
\]
\end{proof}

% ===================================================================
\section{Demostración 2: búsqueda exitosa
  \texorpdfstring{$\Theta(1 + \alpha/2)$}{Theta(1 + alpha/2)}}
% ===================================================================

\begin{theorem}\label{thm:hit}
El número esperado de comparaciones en una búsqueda exitosa es:
\[
  1 + \frac{\alpha}{2} - \frac{1}{2m}
\]
\end{theorem}

\begin{proof}
Supongamos que los $n$ elementos fueron insertados en orden
$k_1, k_2, \ldots, k_n$. Buscamos un elemento elegido uniformemente
al azar entre los $n$ almacenados.

\textbf{Paso 1: comparaciones para encontrar $k_i$.} El elemento
$k_i$ fue el $i$-ésimo insertado. Los elementos insertados
\textbf{después} de~$k_i$ en el \textbf{mismo bucket} se colocan
delante de~$k_i$ (por el prepend). El número de comparaciones es:
\[
  C_i = 1 + \sum_{r=i+1}^{n} Y_{ri}
\]
donde $Y_{ri} = \mathbf{1}[h(k_r) = h(k_i)]$. El $1$ cuenta la
comparación final con~$k_i$.

\textbf{Paso 2: esperanza de $C_i$.} Por SUHA,
$\Pr[h(k_r) = h(k_i)] = 1/m$:
\[
  E[C_i] = 1 + \sum_{r=i+1}^{n} \frac{1}{m} = 1 + \frac{n - i}{m}
\]

\textbf{Paso 3: promedio sobre todos los elementos.}
\begin{align*}
  E[\text{comparaciones}]
  &= \frac{1}{n} \sum_{i=1}^{n} E[C_i]
   = \frac{1}{n} \sum_{i=1}^{n} \left(1 + \frac{n - i}{m}\right) \\
  &= 1 + \frac{1}{nm} \sum_{i=1}^{n} (n - i)
   = 1 + \frac{1}{nm} \cdot \frac{n(n{-}1)}{2}
   = 1 + \frac{n-1}{2m}
\end{align*}

Expresando en función de $\alpha = n/m$:
\[
  1 + \frac{n-1}{2m} = 1 + \frac{\alpha}{2} - \frac{1}{2m}
  \approx 1 + \frac{\alpha}{2}
\]

Incluyendo el costo $O(1)$ de hash y acceso:
$T_{\text{hit}} = \Theta(1 + \alpha/2)$.
\end{proof}

% ===================================================================
\section{Demostración 3: ambos costos son $O(1)$ cuando
  \texorpdfstring{$\alpha = O(1)$}{alpha = O(1)}}
% ===================================================================

\begin{theorem}\label{thm:constant}
Si la tabla se redimensiona para mantener
$\alpha \leq \alpha_{\max}$ (constante fija), entonces búsqueda,
inserción y eliminación son $O(1)$ en esperanza.
\end{theorem}

\begin{proof}
Si $\alpha \leq \alpha_{\max}$ en todo momento:
\begin{itemize}
  \item Búsqueda fallida:
    $\Theta(1 + \alpha) \leq \Theta(1 + \alpha_{\max}) = O(1)$.
  \item Búsqueda exitosa:
    $\Theta(1 + \alpha/2) \leq \Theta(1 + \alpha_{\max}/2) = O(1)$.
  \item Inserción (con verificación de duplicados): dominada por
    búsqueda fallida: $O(1)$.
  \item Eliminación: dominada por búsqueda exitosa: $O(1)$.
\end{itemize}
El costo de redimensionar es $O(n)$ por resize, pero se amortiza a
$O(1)$ por inserción (demostrado en T05).
\end{proof}

\begin{remark}
Esto es lo que hace que las tablas hash sean $O(1)$ en la práctica:
no es que las operaciones sean $O(1)$ para cualquier~$\alpha$, sino
que el redimensionamiento mantiene $\alpha$ acotado.
\end{remark}

% ===================================================================
\section{Demostración 4: la búsqueda exitosa es estrictamente más
  barata que la fallida}
% ===================================================================

\begin{theorem}\label{thm:hit-cheaper}
Para todo $\alpha > 0$:
$T_{\text{hit}} < T_{\text{miss}}$.
\end{theorem}

\begin{proof}
Comparamos los costos totales (hash + recorrido):
\[
  T_{\text{miss}} = 1 + \alpha, \qquad
  T_{\text{hit}} = 1 + \frac{\alpha}{2} - \frac{1}{2m}
\]

$T_{\text{hit}} < T_{\text{miss}}$ equivale a:
\[
  1 + \frac{\alpha}{2} - \frac{1}{2m} < 1 + \alpha
  \iff -\frac{1}{2m} < \frac{\alpha}{2}
\]
que vale siempre que $\alpha > 0$. \checkmark
\end{proof}

\begin{remark}
En una búsqueda exitosa, encontramos el elemento antes de llegar al
final de la lista (en promedio, a mitad de camino). En una fallida,
recorremos toda la lista.
\end{remark}

% ===================================================================
\section{Demostración 5: prepend $\equiv$ append para
  \texorpdfstring{$C_{\text{hit}}$}{C\_hit}}
% ===================================================================

\begin{theorem}\label{thm:prepend-append}
Si los nuevos elementos se insertan al \textbf{final} de la cadena
(append) en vez de al inicio, el costo de búsqueda exitosa es el
mismo:
\[
  E[C_{\text{hit}}^{\text{append}}] = 1 + \frac{n-1}{2m}
\]
\end{theorem}

\begin{proof}
Con inserción al final, el elemento~$k_i$ está precedido por los
elementos insertados \textbf{antes} que~$k_i$ en el mismo bucket:
\[
  C_i^{\text{append}} = 1 + \sum_{r=1}^{i-1} Y_{ri}'
\]
donde $Y_{ri}' = \mathbf{1}[h(k_r) = h(k_i)]$.
\[
  E[C_i^{\text{append}}] = 1 + \frac{i - 1}{m}
\]

Promediando:
\[
  E[C_{\text{hit}}^{\text{append}}]
  = \frac{1}{n}\sum_{i=1}^{n}\left(1 + \frac{i-1}{m}\right)
  = 1 + \frac{1}{nm}\sum_{i=1}^{n}(i{-}1)
  = 1 + \frac{1}{nm} \cdot \frac{n(n{-}1)}{2}
  = 1 + \frac{n-1}{2m}
\]

Idéntico al caso prepend.
\end{proof}

\begin{remark}
La política de inserción (prepend vs append) no afecta el costo
esperado de búsqueda exitosa. En promedio, el elemento buscado está a
mitad de la lista independientemente del orden, ya que buscamos un
elemento uniformemente aleatorio. La clave algebraica es que
$\sum_{i=1}^{n}(n{-}i) = \sum_{i=1}^{n}(i{-}1) = n(n{-}1)/2$.
\end{remark}

% ===================================================================
\section{Resumen de resultados}
% ===================================================================

\begin{table}[h]
\centering
\begin{tabular}{llp{5.5cm}}
\toprule
\textbf{Resultado} & \textbf{Técnica} & \textbf{Clave} \\
\midrule
$C_{\text{miss}} = \alpha$
  & $E[L_{h(k)}] = \alpha$
  & Recorrer lista completa \\
$C_{\text{hit}} = 1 + (n{-}1)/2m$
  & Contar posteriores en mismo bucket
  & $\Pr[h(k_r){=}h(k_i)] = 1/m$, sumar $n{-}i$ \\
$O(1)$ con $\alpha = O(1)$
  & Redimensionamiento
  & $\alpha$ acotado $\Rightarrow$ costos acotados \\
$T_{\text{hit}} < T_{\text{miss}}$ siempre
  & Comparación directa
  & Encontrar antes de fin de lista \\
Prepend $\equiv$ append
  & Simetría de la suma
  & $\sum(n{-}i) = \sum(i{-}1) = n(n{-}1)/2$ \\
\bottomrule
\end{tabular}
\end{table}

\end{document}
